\section{What has been checked, and how to review this draft}\label{sec:verification}

\subsection{Independent exact checks}
The companion script \texttt{tools/verify/verify\_tutorial.py} uses Python's
rational arithmetic and does not depend on a commercial LP solver. It
implements the formulas in this draft as a pedagogical checker, not as the
planned high-performance solver. At each visited basis it independently
assembles $A_{\Bset}$ and solves
\[
 A_{\Bset}z_{\Bset}=b,\qquad
 A_{\Bset}d_{\Bset}=-A_q
\]
by rational Gaussian elimination. It compares every nonbasic candidate's
direction and reduced cost with the forest formulas, checks the exchanged
basis and verifies the final ratio dual.
It also solves the revised simplex multiplier system
$\mat{A}_{\Bset}^T\vecg{\pi}=\vect{g}_{\Bset}$ independently at every visited basis and compares
its solution with the forest multipliers. At each pivot, it solves
$\mat{A}_{\Bset}^T\vecg{\eta}=\vect{e}_r$ for the leaving position and checks every
nonbasic coefficient of the resulting pivot row against the forest direction.

Separately, exhaustive enumeration lists all closures on the small instance,
computes exact maximum ratios and maximal ties, and checks the complete
canonical sequence. For the LP value it enumerates capacity intersections
of pairs of closure points in the weight--profit plane; this independent
oracle does not restrict those pairs to the sequence found by the \FS{}.

\input{../build/tutorial/generated/verification_stats.tex}
The supplied deterministic check suite contains \VerifiedInstances{} instances:
the eight-node example, four explicit edge cases, and 32 seeded small random
DAGs with profits of either sign and positive integer weights. It checks
\VerifiedBases{} bases and \VerifiedDirections{} candidate directions.
The additional revised simplex comparisons check \VerifiedMultipliers{}
multiplier systems and \VerifiedPivotRows{} pivot rows.
The eight-node graph has 36 closures. Its first oracle takes eight pivots,
six at zero step, and its capacity-4 LP value is 7.
Both numerical certificates displayed in the main text are verified by
the checker or its returned forest certificate.

These counts describe a finite review suite, not comprehensive testing of
an optimized implementation, not a statistical experiment, and not evidence
of a runtime advantage. The seed and hashes of the fixture and verification
script are saved with the generated results.

\subsection{Rebuilding the tables and PDF}
From the project root, run
\begin{quote}\ttfamily
python3 tools/build\_tutorial.py
\end{quote}
The command reruns the exact checks, regenerates numerical tables and forest
figures, and compiles this document. The single output PDF is
\begin{quote}\ttfamily\small
build/tutorial/forest\_simplex\_tutorial.pdf
\end{quote}
The generated verification record is
\begin{quote}\ttfamily\small
build/tutorial/generated/verification.json
\end{quote}
All intermediate TeX files remain under \texttt{build/tutorial}. The source
instance, shared notation and bibliography each have one authoritative file.
No original project source is altered or required during this rebuild.

\subsection{Suggested order for mathematical review}
The following checkpoints identify concrete places where an error would
affect the proposed general algorithm:
\begin{enumerate}
\item Check the prerequisite direction in \eqref{eq:lpprec} and the
parent--child sign in \eqref{eq:sign}.
\item Check the converse as well as the forward implication of the basis
characterization, \cref{thm:basis}.
\item Check that the normalization correction in \eqref{eq:direction}
preserves every other nonbasic variable.
\item Check the sign of \eqref{eq:rc-edge}, especially when a child points
to its parent in a zero tree.
\item Check the inclusion of decreasing nonforest slacks in the ratio test.
\item Follow pivots 7 and 8, including the unchanged objective at pivot 8.
\item Check the sign convention in the forest dual certificate and distinguish
it from the capacitated LP dual.
\item Check that merging raw ties yields the maximal, possibly disconnected,
canonical block.
\item Check the stop at nonpositive ratios for the capacity-inequality LP.
\end{enumerate}

\subsection{Work remaining after this review version}
The exact-reference derivation above is accompanied by proofs and finite
checks. The following tasks remain open as project deliverables:
\begin{itemize}
\item external mathematical review and comparison with existing forest,
network-simplex and closure algorithms to determine the actual novelty;
\item a C++ implementation with linear-memory data structures;
\item verified local updates, alternate pricing and warm starts;
\item numerical stress tests beyond the exact small-DAG checker;
\item competitive backend comparisons, parameter tuning and reproducible
runtime/memory experiments;
\item extension beyond the declared positive-weight, single-capacity model.
\end{itemize}
In particular, this document supplies no invented benchmark results and makes
no claim that the eventual optimized solver is faster than parametric flow
or a general LP solver.
